
% GA-PRAC-01
% C3 - Actores del proceso y framework de Ingenieria de Requisitos
% Responsable: Herrera Moran Gilmar Javier
% Requiere: \usepackage{longtable}
%           \usepackage{array}

\section{C3. Actores del proceso y framework de Ingeniería de Requisitos}

\subsection{Marco de referencia}

El framework de Ingeniería de Requisitos (IR) \cite{Pohl2025,RuppPohl2015} se organiza en tres grupos de elementos que se aplican a los cuatro escenarios, independientemente del modelo de proceso elegido:

\begin{itemize}
    \item \textbf{Actividades core:} elicitación, documentación y negociación o resolución de conflictos. Producen y consolidan los requisitos \cite{Pohl2025}.
    \item \textbf{Actividades transversales:} validación y gestión de requisitos. Acompañan a las tres actividades core durante todo el proyecto \cite{Pohl2025,SWEBOK_v4_2024}.
    \item \textbf{Artefactos:} metas, escenarios y requisitos orientados a solución, que representan el problema y la solución con distinto nivel de detalle \cite{Pohl2025}.
\end{itemize}

En cada escenario se identifican primero los \textbf{actores del proceso}, es decir, quienes aportan, deciden, analizan, validan o implementan requisitos. Luego se ubican las cinco actividades del framework y, finalmente, se indica cómo se concretan los tres artefactos. La identificación de actores como partes interesadas y la trazabilidad de los requisitos siguen además las prácticas de la norma ISO/IEC/IEEE 29148 \cite{ISO29148}. Las tablas siguientes usan únicamente los datos de cada escenario y los modelos recomendados en C4.

% ---------------------------------------------------------------------
% ESCENARIO 1
% ---------------------------------------------------------------------
\subsection{Escenario 1: Extractora Río Lindo}

\textbf{Actores y responsabilidades.}

{\small
\setlength{\tabcolsep}{4pt}
\renewcommand{\arraystretch}{1.3}
\begin{longtable}{
|>{\raggedright\arraybackslash}p{0.22\textwidth}
|>{\raggedright\arraybackslash}p{0.16\textwidth}
|>{\raggedright\arraybackslash}p{0.50\textwidth}|}
\hline
\textbf{Actor} & \textbf{Tipo} & \textbf{Responsabilidad en el proceso de IR} \\
\hline
\endfirsthead
\hline
\textbf{Actor} & \textbf{Tipo} & \textbf{Responsabilidad en el proceso de IR} \\
\hline
\endhead
Jefe de planta & Stakeholder / decisor & Aporta las reglas operativas del proceso, aprueba requisitos y prioriza. Está disponible solo dos mañanas por semana, por lo que sus sesiones deben planificarse. \\
\hline
Responsable de seguridad & Stakeholder / aprobador & Define y aprueba los requisitos de protección (alarmas, bloqueos, límites de temperatura y presión). \\
\hline
Operadores de turnos rotativos & Usuarios finales & Describen el uso real de la interfaz y de las alarmas. Validan escenarios de operación en cada turno. \\
\hline
Aseguradora & Stakeholder externo & Exige un documento de requisitos de seguridad aprobado y evidencias de pruebas antes de la puesta en marcha. \\
\hline
Fabricante de los esterilizadores & Fuente de requisitos & Establece los límites técnicos y de seguridad del equipo que el sistema debe respetar. \\
\hline
Proveedores de sensores y válvulas & Stakeholder técnico & Su selección (tercer mes) determina interfaces y restricciones de integración. \\
\hline
Analista de requisitos & Equipo del proyecto & Conduce la elicitación, redacta la especificación, gestiona conflictos y mantiene la trazabilidad. \\
\hline
Equipo de desarrollo e integración & Equipo del proyecto & Implementa el sistema embebido, estima la factibilidad y verifica los requisitos en cada incremento. \\
\hline
\end{longtable}
}

\textbf{Actividades del framework ubicadas en el escenario.}

{\small
\setlength{\tabcolsep}{4pt}
\renewcommand{\arraystretch}{1.3}
\begin{longtable}{
|>{\raggedright\arraybackslash}p{0.20\textwidth}
|>{\raggedright\arraybackslash}p{0.13\textwidth}
|>{\raggedright\arraybackslash}p{0.55\textwidth}|}
\hline
\textbf{Actividad} & \textbf{Tipo} & \textbf{Aplicación en el escenario} \\
\hline
\endfirsthead
\hline
\textbf{Actividad} & \textbf{Tipo} & \textbf{Aplicación en el escenario} \\
\hline
\endhead
Elicitación & Core & Entrevistas con el jefe de planta en sus dos mañanas semanales, observación de operadores en distintos turnos y análisis de las especificaciones del fabricante y de la aseguradora. \\
\hline
Documentación & Core & Especificación formal de los requisitos de seguridad (umbrales, bloqueos y alarmas), seguida de la documentación incremental de interfaz, reportes y comunicación con sensores. \\
\hline
Negociación y resolución de conflictos & Core & Equilibrar el deseo de operadores de una interfaz cómoda con las restricciones de seguridad, y los plazos del proyecto con la incertidumbre de los sensores. \\
\hline
Validación & Transversal & Aprobación formal de los requisitos de seguridad por el responsable de seguridad y la aseguradora. Revisión de prototipos de interfaz con operadores de varios turnos al cierre de cada incremento. \\
\hline
Gestión de requisitos & Transversal & Línea base congelada para las reglas de seguridad, trazabilidad entre requisito, prueba y evidencia, y control formal de cambios cuando se definan los sensores. \\
\hline
\end{longtable}
}

\textbf{Artefactos del framework.}

{\small
\setlength{\tabcolsep}{4pt}
\renewcommand{\arraystretch}{1.3}
\begin{longtable}{
|>{\raggedright\arraybackslash}p{0.20\textwidth}
|>{\raggedright\arraybackslash}p{0.68\textwidth}|}
\hline
\textbf{Artefacto} & \textbf{Contenido en el escenario} \\
\hline
\endfirsthead
\hline
\textbf{Artefacto} & \textbf{Contenido en el escenario} \\
\hline
\endhead
Metas & Proteger a los operadores y la producción; controlar temperatura, presión y niveles de tanques dentro de los límites del fabricante; cumplir las exigencias de la aseguradora. \\
\hline
Escenarios & Un operador recibe una alarma por presión elevada y el sistema activa el bloqueo. Un operador de otro turno consulta el reporte del lote. Se pierde la lectura de un sensor y el sistema pasa a estado seguro. \\
\hline
Requisitos orientados a solución & Requisitos funcionales de lectura, alarma y bloqueo; requisitos de calidad como tiempo de respuesta y disponibilidad; restricciones de integración con los sensores y válvulas que se seleccionen. \\
\hline
\end{longtable}
}

% ---------------------------------------------------------------------
% ESCENARIO 2
% ---------------------------------------------------------------------
\subsection{Escenario 2: Sabores de la Abuela Rosa}

\textbf{Actores y responsabilidades.}

{\small
\setlength{\tabcolsep}{4pt}
\renewcommand{\arraystretch}{1.3}
\begin{longtable}{
|>{\raggedright\arraybackslash}p{0.22\textwidth}
|>{\raggedright\arraybackslash}p{0.16\textwidth}
|>{\raggedright\arraybackslash}p{0.50\textwidth}|}
\hline
\textbf{Actor} & \textbf{Tipo} & \textbf{Responsabilidad en el proceso de IR} \\
\hline
\endfirsthead
\hline
\textbf{Actor} & \textbf{Tipo} & \textbf{Responsabilidad en el proceso de IR} \\
\hline
\endhead
Doña Rosa (propietaria) & Cliente / decisora & Aporta las necesidades del negocio, prioriza el contenido de cada ciclo y valida los incrementos. Está disponible a diario. \\
\hline
Hijo de la propietaria & Stakeholder & Propone ideas y mejoras. Sus solicitudes deben pasar por la priorización de Doña Rosa para evitar ampliar el alcance. \\
\hline
Clientes del emprendimiento & Usuarios finales & Realizan pedidos y pagos por transferencia desde teléfonos Android de gama baja. Sus datos personales deben protegerse. \\
\hline
Personal de cocina o entrega & Usuarios internos & Reciben y actualizan el estado de los pedidos. \\
\hline
Dos desarrolladores a medio tiempo & Equipo del proyecto & Asumen elicitación, análisis, implementación y pruebas, con capacidad limitada. Uno de ellos actúa como analista de requisitos. \\
\hline
Autoridad de protección de datos (LOPDP) & Marco normativo & Establece las obligaciones sobre el tratamiento de datos personales de los clientes. \\
\hline
\end{longtable}
}

\textbf{Actividades del framework ubicadas en el escenario.}

{\small
\setlength{\tabcolsep}{4pt}
\renewcommand{\arraystretch}{1.3}
\begin{longtable}{
|>{\raggedright\arraybackslash}p{0.20\textwidth}
|>{\raggedright\arraybackslash}p{0.13\textwidth}
|>{\raggedright\arraybackslash}p{0.55\textwidth}|}
\hline
\textbf{Actividad} & \textbf{Tipo} & \textbf{Aplicación en el escenario} \\
\hline
\endfirsthead
\hline
\textbf{Actividad} & \textbf{Tipo} & \textbf{Aplicación en el escenario} \\
\hline
\endhead
Elicitación & Core & Conversaciones diarias breves con Doña Rosa, observación de cómo toma pedidos hoy, y consulta a algunos clientes sobre cómo piden y pagan. \\
\hline
Documentación & Core & Historias de usuario con criterios de aceptación en un \textit{backlog} vivo. Documentación breve y obligatoria de privacidad de datos y de conciliación de transferencias. \\
\hline
Negociación y resolución de conflictos & Core & Decidir qué ideas del hijo entran en cada ciclo y cuáles se aplazan, según presupuesto y plazo de 14 semanas. Resolver el choque entre agilidad comercial y cumplimiento legal. \\
\hline
Validación & Transversal & Doña Rosa prueba cada incremento en teléfonos de gama baja y con conexión intermitente. Se verifica la protección de datos y el control de pagos antes de cerrar cada ciclo. \\
\hline
Gestión de requisitos & Transversal & Priorización continua (por ejemplo, MoSCoW), control del alcance de cada ciclo y registro de cambios del menú y de las promociones. \\
\hline
\end{longtable}
}

\textbf{Artefactos del framework.}

{\small
\setlength{\tabcolsep}{4pt}
\renewcommand{\arraystretch}{1.3}
\begin{longtable}{
|>{\raggedright\arraybackslash}p{0.20\textwidth}
|>{\raggedright\arraybackslash}p{0.68\textwidth}|}
\hline
\textbf{Artefacto} & \textbf{Contenido en el escenario} \\
\hline
\endfirsthead
\hline
\textbf{Artefacto} & \textbf{Contenido en el escenario} \\
\hline
\endhead
Metas & Recibir y gestionar pedidos de comida típica de forma ordenada; controlar los pagos por transferencia; proteger los datos de los clientes; tener la aplicación lista para las fiestas dentro de 14 semanas y del presupuesto. \\
\hline
Escenarios & Un cliente elige un plato del menú del día y confirma el pedido. Un cliente con mala conexión intenta pagar y el sistema conserva el pedido. Doña Rosa cambia el menú y las promociones. \\
\hline
Requisitos orientados a solución & Funcionales: catálogo editable, carrito, registro de pedidos y de transferencias. De calidad: funcionamiento en Android de gama baja y con conectividad intermitente. Restricciones: protección de datos personales y presupuesto cerrado. \\
\hline
\end{longtable}
}

% ---------------------------------------------------------------------
% ESCENARIO 3
% ---------------------------------------------------------------------
\subsection{Escenario 3: Sistema de trámites municipales}

\textbf{Actores y responsabilidades.}

{\small
\setlength{\tabcolsep}{4pt}
\renewcommand{\arraystretch}{1.3}
\begin{longtable}{
|>{\raggedright\arraybackslash}p{0.22\textwidth}
|>{\raggedright\arraybackslash}p{0.16\textwidth}
|>{\raggedright\arraybackslash}p{0.50\textwidth}|}
\hline
\textbf{Actor} & \textbf{Tipo} & \textbf{Responsabilidad en el proceso de IR} \\
\hline
\endfirsthead
\hline
\textbf{Actor} & \textbf{Tipo} & \textbf{Responsabilidad en el proceso de IR} \\
\hline
\endhead
Administrador del contrato (GAD) & Cliente / aprobador & Aprueba los entregables, firma las actas y coordina con el municipio. Tiene disponibilidad semanal. \\
\hline
Funcionarios de ventanilla y rentas & Usuarios finales & Describen los flujos de permisos, patentes y certificados, y los cobros. Participan de forma irregular. \\
\hline
Ciudadanos solicitantes & Usuarios finales & Realizan y consultan trámites, y entregan datos personales que deben protegerse. \\
\hline
Personal de tecnología municipal & Stakeholder técnico & Provee acceso a los servidores y conocimiento de la base de datos antigua sin documentación. \\
\hline
Autoridades municipales (ordenanza de tasas) & Stakeholder normativo & La posible ordenanza del cuarto mes puede cambiar tasas y flujos de aprobación. \\
\hline
Entes de contratación y control & Marco contractual & Los pliegos fijan los tres entregables, las fechas, los pagos y las multas por retraso. \\
\hline
Contratista: líder y analista de requisitos & Equipo del proyecto & Elicita, redacta la especificación aprobada en el mes 2, gestiona cambios y mantiene la trazabilidad. \\
\hline
Contratista: desarrollo y pruebas (seis personas en total) & Equipo del proyecto & Implementa e integra los módulos, investiga la base de datos antigua y genera evidencias de pruebas. \\
\hline
\end{longtable}
}

\textbf{Actividades del framework ubicadas en el escenario.}

{\small
\setlength{\tabcolsep}{4pt}
\renewcommand{\arraystretch}{1.3}
\begin{longtable}{
|>{\raggedright\arraybackslash}p{0.20\textwidth}
|>{\raggedright\arraybackslash}p{0.13\textwidth}
|>{\raggedright\arraybackslash}p{0.55\textwidth}|}
\hline
\textbf{Actividad} & \textbf{Tipo} & \textbf{Aplicación en el escenario} \\
\hline
\endfirsthead
\hline
\textbf{Actividad} & \textbf{Tipo} & \textbf{Aplicación en el escenario} \\
\hline
\endhead
Elicitación & Core & Entrevistas con el administrador del contrato, observación del trabajo en ventanilla y rentas, análisis de los pliegos y de la ordenanza vigente, e ingeniería inversa de la base de datos antigua. \\
\hline
Documentación & Core & Especificación de requisitos formal para el mes 2, con trazabilidad, actas de aceptación y registros de pruebas necesarios para justificar los entregables y los pagos. \\
\hline
Negociación y resolución de conflictos & Core & Acordar con el municipio qué cambios (flujo de aprobación, nueva ordenanza) caben dentro del alcance contratado y cuáles requieren una modificación contractual. \\
\hline
Validación & Transversal & Aceptación formal de la especificación en el mes 2, validación periódica de módulos con funcionarios, pruebas de integración y acta de entrega-recepción en el mes 7. \\
\hline
Gestión de requisitos & Transversal & Línea base contractual, control formal de cambios con evaluación de impacto, versionado de la especificación y trazabilidad entre requisitos, pruebas y entregables. \\
\hline
\end{longtable}
}

\textbf{Artefactos del framework.}

{\small
\setlength{\tabcolsep}{4pt}
\renewcommand{\arraystretch}{1.3}
\begin{longtable}{
|>{\raggedright\arraybackslash}p{0.20\textwidth}
|>{\raggedright\arraybackslash}p{0.68\textwidth}|}
\hline
\textbf{Artefacto} & \textbf{Contenido en el escenario} \\
\hline
\endfirsthead
\hline
\textbf{Artefacto} & \textbf{Contenido en el escenario} \\
\hline
\endhead
Metas & Gestionar cerca de 9\,000 trámites anuales; cumplir los tres hitos contractuales sin multas; proteger los datos personales; integrar el sistema con los servidores y la base de datos municipales. \\
\hline
Escenarios & Un ciudadano solicita un permiso de funcionamiento. Un funcionario de rentas cobra una patente con las tasas vigentes. Se aprueba la nueva ordenanza y las tasas cambian. Se emite un certificado de no adeudar consultando la base antigua. \\
\hline
Requisitos orientados a solución & Funcionales: registro, flujo de aprobación, cobro y emisión de documentos. De calidad: seguridad y protección de datos. Restricciones: integración con la base antigua y servidores municipales, y entregables contractuales. \\
\hline
\end{longtable}
}

% ---------------------------------------------------------------------
% ESCENARIO 4
% ---------------------------------------------------------------------
\subsection{Escenario 4: Cooperativa Frutos del Estero}

\textbf{Actores y responsabilidades.}

{\small
\setlength{\tabcolsep}{4pt}
\renewcommand{\arraystretch}{1.3}
\begin{longtable}{
|>{\raggedright\arraybackslash}p{0.22\textwidth}
|>{\raggedright\arraybackslash}p{0.16\textwidth}
|>{\raggedright\arraybackslash}p{0.50\textwidth}|}
\hline
\textbf{Actor} & \textbf{Tipo} & \textbf{Responsabilidad en el proceso de IR} \\
\hline
\endfirsthead
\hline
\textbf{Actor} & \textbf{Tipo} & \textbf{Responsabilidad en el proceso de IR} \\
\hline
\endhead
Gerente de la cooperativa & Cliente / enlace & Disponible casi a diario. Aporta las reglas de negocio y valida funcionalidades, pero no aprueba cambios importantes por sí sola. \\
\hline
Directiva & Decisor & Se reúne mensualmente y aprueba los cambios importantes de alcance. \\
\hline
Socios productores (640) & Usuarios finales & Consultan entregas, pagos, créditos y saldos. Participan principalmente en días de acopio y en asambleas trimestrales. \\
\hline
Personal de acopio & Usuarios internos & Registran peso, humedad y calidad de las entregas, incluso sin conexión. \\
\hline
Compradores exportadores & Stakeholder externo & Sus exigencias de trazabilidad y certificación del cacao pueden cambiar. \\
\hline
Donante & Financiador & Exige informes semestrales de avance, por lo que condiciona los hitos y las evidencias. \\
\hline
Organismo supervisor de la economía popular y solidaria & Marco normativo & Control institucional y cumplimiento del estatuto de la cooperativa. \\
\hline
Equipo de cuatro personas & Equipo del proyecto & Elicita en campo, especifica los cálculos de liquidación, implementa el registro sin conexión y la sincronización, y prepara los informes. \\
\hline
\end{longtable}
}

\textbf{Actividades del framework ubicadas en el escenario.}

{\small
\setlength{\tabcolsep}{4pt}
\renewcommand{\arraystretch}{1.3}
\begin{longtable}{
|>{\raggedright\arraybackslash}p{0.20\textwidth}
|>{\raggedright\arraybackslash}p{0.13\textwidth}
|>{\raggedright\arraybackslash}p{0.55\textwidth}|}
\hline
\textbf{Actividad} & \textbf{Tipo} & \textbf{Aplicación en el escenario} \\
\hline
\endfirsthead
\hline
\textbf{Actividad} & \textbf{Tipo} & \textbf{Aplicación en el escenario} \\
\hline
\endhead
Elicitación & Core & Entrevistas con la gerente y la directiva, observación del proceso actual (tarjetas de papel y hojas de cálculo) en los días de acopio, y consulta a socios en asambleas y recintos. \\
\hline
Documentación & Core & Especificación detallada de reglas de liquidación, créditos y saldos; requisitos de funcionamiento sin conexión y sincronización; documentación de trazabilidad y evidencias para los informes semestrales. \\
\hline
Negociación y resolución de conflictos & Core & Conciliar las prioridades de la gerente, de la directiva y de los socios, y los cambios que piden los compradores con el plazo de doce meses y la aprobación mensual de la directiva. \\
\hline
Validación & Transversal & Validación periódica de incrementos con la gerente, con socios durante los días de acopio y en asambleas trimestrales, y aprobación de la directiva en los cambios importantes. \\
\hline
Gestión de requisitos & Transversal & Línea base para las reglas de liquidación (baja volatilidad), control de cambios de trazabilidad y certificación, y vínculo de cada incremento con los informes del donante. \\
\hline
\end{longtable}
}

\textbf{Artefactos del framework.}

{\small
\setlength{\tabcolsep}{4pt}
\renewcommand{\arraystretch}{1.3}
\begin{longtable}{
|>{\raggedright\arraybackslash}p{0.20\textwidth}
|>{\raggedright\arraybackslash}p{0.68\textwidth}|}
\hline
\textbf{Artefacto} & \textbf{Contenido en el escenario} \\
\hline
\endfirsthead
\hline
\textbf{Artefacto} & \textbf{Contenido en el escenario} \\
\hline
\endhead
Metas & Registrar las entregas de cacao con exactitud; calcular liquidaciones correctas; permitir a los socios consultar entregas, pagos, créditos y saldos; cumplir con el donante y el organismo de control. \\
\hline
Escenarios & Un encargado registra una entrega sin conexión en un recinto y la sincroniza después. Un socio consulta su saldo. La gerente calcula una liquidación. Un comprador exige un nuevo dato de trazabilidad. \\
\hline
Requisitos orientados a solución & Funcionales: registro de entregas, liquidaciones, créditos y consultas. De calidad: exactitud de cálculos, integridad de los datos y protección de datos personales. Restricciones: trabajo sin conexión, sincronización posterior y estatuto de la cooperativa. \\
\hline
\end{longtable}
}

\subsection{Síntesis}

En los cuatro escenarios, el análisis de los actores muestra dónde se concentran las decisiones y qué limita la participación: el jefe de planta y la aseguradora en la Extractora Río Lindo, la propietaria en Sabores de la Abuela Rosa, el administrador del contrato en el GAD Municipal y la directiva junto con el donante en la Cooperativa Frutos del Estero.

Las cinco actividades del framework son las mismas en todos los casos, pero cambia su peso. Cuando predominan la criticidad y las obligaciones formales (Escenarios 1 y 3), destacan la documentación y la gestión de requisitos. Cuando predominan el cambio y la disponibilidad del cliente (Escenario 2), destacan la elicitación continua y la negociación del alcance. En el Escenario 4 la validación debe adaptarse a los días de acopio y a las asambleas. Esta lectura es coherente con los modelos recomendados en C4.
